\section{Scripts}
\label{sec:extension-scripts-scripts}

\begin{figure}[htb!]
	\centering
	\includegraphics[width=0.6\textwidth]{exemple-contrainte-regles-gauthier}
	\caption{Exemple de décousure d'un volume par Gauthier}
	\label{fig:reevaluation-scripts-scripts-scripts-jerboa-non-decousure}
\end{figure}

\begin{figure}[htb!]
	\centering
	\includegraphics[height=40pt]{rule-unsew-volume-a3}
	\caption{Règle de décousure d'un volume par ses \(3\)-liaisons}
	\label{fig:reevaluation-scripts-scripts-scripts-jerboa-regle-decousure-volume}
\end{figure}

\begin{figure}[htb!]
	\centering
	\includegraphics[height=40pt]{rule-unsew-face-a3}
	\caption{Règle de décousure d'une face par ses \(3\)-liaisons}
	\label{fig:reevaluation-scripts-scripts-scripts-jerboa-regle-decousure-face}
\end{figure}

Les règles Jerboa nous permettent de représenter diverses opérations de
modélisations plus ou moins complexes. %
Toutefois, comme nous l'avons vu dans le
chapitre~\ref{cha:outils-modelisation-base-regles} avec la
définition~\ref{def:orbite-instanciation}
(page~\pageref{def:orbite-instanciation}), les orbites filtrées par la partie
gauche d'une règle doivent être isomorphes à renommage près. %
%% Exemple de problème lié à l'isomorphisme des règles Jerboa
La figure~\ref{fig:reevaluation-scripts-scripts-scripts-jerboa-non-decousure}
illustre un exemple où la décousure d'un volume par ses 3-liaisons
(figure~\ref{fig:reevaluation-scripts-scripts-scripts-jerboa-regle-decousure-volume})
est impossible, car certaines des faces du volume sont libres. %
Elles ne possèdent pas de 3-liaisons à découdre et le volume
\textcolor{orange}{orange} ne peut pas être déconnecté du reste de la composante
connexe. %
Pour déconnecter ce volume, il faudrait appliquer la règle de décousure de face
(figure~\ref{fig:reevaluation-scripts-scripts-scripts-jerboa-regle-decousure-face})
sur chacune des faces incidentes à un autre volume. %



\subsection{Scripts Jerboa}
\label{sec:extension-scripts-scripts-scripts-jerboa}



Les scripts Jerboa~\cite{belhaouari2014jerboa,gauth2019} permettent, à l'aide
d'un langage de script, de construire des opérations avec des comportements plus
complexes que les règles. %
Un script peut faciliter, par exemple, le processus de décousure du volume
\textcolor{orange}{orange} en évitant de réappliquer manuellement la règle de
décousure de face
(figure~\ref{fig:reevaluation-scripts-scripts-scripts-jerboa-non-decousure}). %
Le langage utilisé pour écrire les scripts Jerboa définit l'ensemble des
instructions nécessaires pour manipuler les G-cartes et leurs orbites. %
Les structures de contrôles classiques de Java et du C++ sont disponibles, à
savoir les boucles (\texttt{for}, \texttt{foreach}, \texttt{while}) et les
alternatives. %
Plus particulièrement, ces scripts peuvent appeler d'autres opérations. %
Ces appels sont les instructions élémentaires d'un script qui, dans sa forme la
plus simple, est une séquence d'appels de règles. %

%% exemple de script pour résoudre le problème
\begin{figure}[htb!]
	\centering
	\includegraphics[height=40pt]{script-unsew012-a3}
	\caption{Motif gauche du script de décousure de volume}
	\label{fig:reevaluation-scripts-scripts-scripts-jerboa-graphe-script}
\end{figure}

\begin{lstlisting}[style=script-style,language=jerboa-script,caption=Script de
décousure d'un volume,
label=lst:reevaluation-scripts-scripts-scripts-jerboa-script-decousure-volume]
for(JerboaDart face : <0,1,2>_<0,1>(volume#0))
{
    try
    {
        @rule<UnsewFace01A3>(face);
    } catch (JerboaException e){
        // echec de l'application
    }
}
return null;
\end{lstlisting}

Les scripts, comme les règles, sont des opérations. %
À ce titre, elles disposent de paramètres en entrée qui peuvent être des
paramètres topologiques ou de plongement (géométrie, couleur, etc.). %
En pratique, ces paramètres sont utilisés différemment. %
%% Les entrées et les sorties : montrer avec le script
Observons le motif gauche du script qui découd un volume de sa composante
connexe
(figure~\ref{fig:reevaluation-scripts-scripts-scripts-jerboa-graphe-script}). %
Ce motif est composé d'un unique n\oe ud d'accroche appelé \texttt{volume}. %
Son étiquette désigne un brin (\(\orbit{}\)) ainsi qu'une information de
\textit{multiplicité} (\texttt{{1}}) qui indique le nombre d'orbites sur
lesquelles ce script peut être appliqué, ici une seule orbite. %
%% Les paramètres topologiques : montrer avec le script
Une différence notable, entre les motifs gauches des règles et de scripts, est
que les motifs de scripts ne filtrent pas d'orbites. %
Ceux-ci récupèrent simplement une liste de brin, un par orbite, et délèguent le
filtrage aux règles qu'ils appellent. %
Étudions maintenant le script
(liste~\ref{lst:reevaluation-scripts-scripts-scripts-jerboa-script-decousure-volume}). %
%% Structures dans un script
La principale structure de ce script est une boucle qui itère chacune des faces
d'un volume incident à un brin obtenu par le n\oe ud \texttt{volume}. %
L'itération est formalisée par l'instruction
\texttt{<0,1,3>\_<0,1>(volume\#{}0)} qui retourne une liste de brins à raison
d'un brin par face. %
Notons, d'une part que \texttt{volume\#{}0} est le premier (et unique brin)
contenu dans \texttt{volume} et, d'autre part, que la variable \texttt{face} est
effectivement un brin de type \texttt{JerboaDart}. %
%% Appels de scripts dans les scripts
Dans cette boucle, un bloc \texttt{try-catch} contient un appel de la règle
\texttt{UnsewFace01A3} qui doit alors être appliquée à chacune des faces d'un
cube. %
Celui-ci empêche une sortie prématurée du programme lorsque la règle échoue
suite à son application sur une face libre. %
Enfin, et contrairement aux règles pour lesquelles le motif droit construit la
sortie d'une application, une telle sortie doit être explicitement construite
dans le corps d'un script. %
\maxime{Il manque peut-être encore un morceau, pour expliquer ce qu'est une sortie de script, comment ça se construit et avec quoi (\texttt{JerboaRuleResult}). Y revenir pour la correction. }

Ainsi, les scripts Jerboa permettent d'écrire des opérations plus complexes que
les règles puisqu'ils proposent d'organiser des opérations à l'aide de
structures de contrôle. %
Une opération, voire une séquence d'opérations, peut être répété un certain
nombre de fois à l'aide de boucles. %
Une application peut être soumise à condition et le script peut contenir une
alternative selon l'évaluation de la condition. %
Enfin, puisqu'une règle et un script sont tous deux des opérations, un script
peut adopter un comportement d'autant plus complexe qu'il peut appeler un autre
script. %


\subsection{Scripts de réévaluation}
\label{sec:extension-scripts-scripts-scripts-reevaluation}


Dans ce chapitre, les scripts que nous proposons de réévaluer dans ce chapitre
comportent qu'un sous-ensemble d'instruction des scripts Jerboa. %
Nos travaux étant encore à un stade préliminaire, nous proposons de traiter que
trois structures de contrôle pour les instructions. %

\begin{description}
  \item[Séquence]

        ~

\begin{lstlisting}[style=script-style,language=jerboa-script]
@rule<ruleA>(brin);
@rule<ruleB>(brin);
@rule<ruleC>(brin);
\end{lstlisting}

  \item[Boucle «~foreach~»]

        ~

\begin{lstlisting}[style=script-style,language=jerboa-script]
for(JerboaDart face : <0,1,2>_<0,1>(volume#0))
{
    @rule<ruleA>(face);
}
\end{lstlisting}

  \item[Alternative]

        ~

\begin{lstlisting}[style=script-style,language=jerboa-script]
if(CONDITION)
{
    @rule<ruleA>(brin);
} else {
    @rule<ruleB>(brin);
}
\end{lstlisting}
\end{description}


%%% Local Variables:
%%% mode: LaTeX
%%% TeX-master: "../../../main"
%%% End:
